\section{The exact passage from finite sets to Kronecker multiplicities}

\subsection*{Finding M-001: a universal signed orbit formula}
Fix \(N\geq 0\).  For a weak composition
\(\alpha=(\alpha_1,\ldots,\alpha_l)\) of \(N\), let
\(X_\alpha\) be the \(S_N\)-set of ordered decompositions
\([N]=A_1\sqcup\cdots\sqcup A_l\), with \(|A_i|=\alpha_i\).
If a coordinate of \(\alpha\) is negative, put \([X_\alpha]=0\).
For three weak compositions of \(N\), write
\[
 C(\alpha,\beta,\gamma)=
 \#\left\{(u_{ijk})\in\mathbb Z_{\geq0}^{l\times m\times r}:
 \sum_{j,k}u_{ijk}=\alpha_i,\ 
 \sum_{i,k}u_{ijk}=\beta_j,\ 
 \sum_{i,j}u_{ijk}=\gamma_k\right\}.
\]
The same convention makes \(C=0\) if any margin is negative.
For partitions \(\lambda,\mu,\nu\) of \(N\), of lengths \(l,m,r\),
respectively, define
\(\alpha(\pi)_i=\lambda_i-i+\pi(i)\), and similarly
\(\beta(\sigma)\), \(\gamma(\tau)\).  Then
\begin{equation}\label{eq:M001}
 g(\lambda,\mu,\nu)=
 \sum_{\pi\in S_l,\,\sigma\in S_m,\,\tau\in S_r}
 \operatorname{sgn}(\pi)\operatorname{sgn}(\sigma)
 \operatorname{sgn}(\tau)
 C\bigl(\alpha(\pi),\beta(\sigma),\gamma(\tau)\bigr).
\end{equation}
This is a derived form of the Jacobi--Trudi identity, not a claim
of a new formula in the literature.

\begin{proof}
Let \(\ell:B(S_N)\longrightarrow R_{\mathbb Q}(S_N)\) be linearization.
The Young permutation module \(\mathbb Q[X_\alpha]\) has Frobenius
characteristic \(h_{\alpha_1}\cdots h_{\alpha_l}\).
Jacobi--Trudi therefore gives
\[
 [S^\lambda]=
 \ell\left(\sum_{\pi\in S_l}
                  \operatorname{sgn}(\pi)[X_{\alpha(\pi)}]\right).
\]
The precise classical inputs are Macdonald, \emph{Symmetric Functions
and Hall Polynomials}, second edition, I, (3.4) and (3.4'), and I, \S7, equations (7.2)--(7.6);
the Frobenius map and the tensor-product inner product are also stated
in Rosas, \S2, equations (3)--(5).
All irreducible characters of \(S_N\) are real and self-dual.
Consequently \(g(\lambda,\mu,\nu)\) is the dimension of the
\(S_N\)-invariants in \(S^\lambda\otimes S^\mu\otimes S^\nu\).
The dimension of the invariants of a permutation module
\(\mathbb Q[Y]\) is the number of orbits of \(Y\).
On substituting the three displayed virtual expressions, it remains
to determine the orbits of
\(X_\alpha\times X_\beta\times X_\gamma\).
The integers \(u_{ijk}=|A_i\cap B_j\cap C_k|\) are orbit invariants.
Conversely, given two triples with the same integers, choose a
bijection separately on each corresponding triple intersection;
their disjoint union is a permutation carrying one triple to the
other.  Every nonnegative table with the required margins is realized
by partitioning \([N]\) into cells of the prescribed sizes.
Thus the orbit number is exactly \(C(\alpha,\beta,\gamma)\), proving
\eqref{eq:M001}.
\end{proof}

\paragraph{What the formula does and does not provide.}
It recovers every Kronecker coefficient from actual finite-set orbit
counts and integral changes of basis.  The determinant signs are
essential to this argument.  It is not a cancellation-free count,
and it supplies neither saturation nor a polynomial-time algorithm
for binary-encoded arbitrary partitions.  The effective-lift
obstruction in A-003 explains why simply replacing the virtual
Young expression by an effective lift cannot work.

\subsection*{Finding M-003: the rational gamma graded ring loses this information}
Let \(G\) be any finite group, let \(R=R_{\mathbb C}(G)\) have its
exterior-power special lambda structure and rank augmentation, and
let \(I=\ker(\dim:R\to\mathbb Z)\).
For the gamma filtration of SGA~6, Exp.~V, 3.10,
\begin{equation}\label{eq:M003}
 F_\gamma^rR\otimes\mathbb Q=I\otimes\mathbb Q
 \quad(r\geq1),\qquad
 \operatorname{gr}_\gamma^rR\otimes\mathbb Q=0
 \quad(r\geq1).
\end{equation}
This statement concerns rationalization of this filtration, not its
integral associated graded ring.

\begin{proof}
In \(R_{\mathbb Q}:=R\otimes_{\mathbb Z}\mathbb Q\) put
\(e=[\mathbb C[G]]/|G|\).
The character of \(e\) is one at the identity and zero elsewhere.
Characters determine representation-ring classes, so \(e^2=e\) and
\(xe=\dim(x)e\) for every \(x\in R_{\mathbb Q}\).
It follows that \(I_{\mathbb Q}=(1-e)R_{\mathbb Q}\), and hence
\(I_{\mathbb Q}^r=I_{\mathbb Q}\) for \(r\geq1\).
By definition \(F_\gamma^rR\) is generated by products
\(\prod_j\gamma^{i_j}(x_j)\), where
\(x_j\in I\), \(i_j>0\), and \(\sum_j i_j\geq r\).
Since \(\gamma^1(x)=x\), it contains \(I^r\).
The augmentation is a lambda homomorphism; for \(x\in I\),
\(\dim(\gamma_t(x))=1\).  Thus every positive gamma operation on
\(x\) has rank zero and \(F_\gamma^rR\subset I\).
Tensoring \(I^r\subset F_\gamma^rR\subset I\) with the flat
\(\mathbb Z\)-module \(\mathbb Q\) proves the first equality.
Exactness of rationalization proves the assertion about the
associated graded.
\end{proof}

\paragraph{Consequence.}
The degree-zero rational graded piece retains the augmentation,
namely dimension; every positive graded piece is zero. Passing
to this associated graded cannot retain the additional representation
data needed for shape-dependent Kronecker multiplicities. In particular, a
saturation argument that uses only that associated graded has
discarded the data it would need.  This is a precise loss-of-data
statement; it does not rule out arguments using the full lambda
ring or integral gamma information.

\subsection*{Finding M-004: an integral invariant scheme, with its good-prime qualification}
Assume \(N\geq1\). Choose stable full integral lattices in the rational Specht modules, denoted \(S^\lambda_{\mathbb Z}\),
\(S^\mu_{\mathbb Z}\), \(S^\nu_{\mathbb Z}\), and put
\[
 T=S^\lambda_{\mathbb Z}\otimes S^\mu_{\mathbb Z}
                  \otimes S^\nu_{\mathbb Z}.
\]
Let \(s_1,\ldots,s_{N-1}\) be the adjacent transpositions of \(S_N\).
The integral map
\[
 A:T\longrightarrow T^{\,N-1},\qquad
 v\longmapsto((s_i-1)v)_i
\]
has
\begin{equation}\label{eq:M004}
 \operatorname{rank}_{\mathbb Z}\ker A
       =g(\lambda,\mu,\nu).
\end{equation}
Moreover the affine group scheme of solutions \(Av=0\) over
\(\mathbb Z[1/N!]\) is a vector bundle of this rank.  Formation of
its module of sections commutes with arbitrary base change within
that open subscheme of \(\operatorname{Spec}\mathbb Z\).

\begin{proof}
The equations \(s_iv=v\) for generators are equivalent to invariance
under \(S_N\).  Such lattices exist: take the integer span of all group translates of a rational basis. They are finite free, and their rationalizations are the corresponding irreducible rational representations.  Since localization is exact,
\((\ker A)\otimes\mathbb Q=(T\otimes\mathbb Q)^{S_N}\).
The latter has dimension \(g\), by self-duality and the definition
of a tensor-product multiplicity.  A submodule of a finite free
\(\mathbb Z\)-module is free, proving \eqref{eq:M004}.
Over \(B=\mathbb Z[1/N!]\), the averaging operator
\[
 P=\frac{1}{N!}\sum_{\sigma\in S_N}\sigma
\]
is an idempotent with image precisely \(T_B^{S_N}\).
This image is a direct summand of a finite free module.  It is
therefore finite projective, and its rank is constant on the
connected scheme \(\operatorname{Spec}B\); the generic rank is \(g\).
The same idempotent exists after any \(B\)-algebra base change, so
its image commutes with that base change.  Finally the equations
\(Av=0\) represent the functor of fixed vectors, and over \(B\)
this functor is the vector bundle associated to the indicated
projective module.
\end{proof}

\paragraph{Limits of the model.}
At primes dividing \(N!\), invariant formation need not commute
with base change, so the same dimension assertion has not been
made there.  The constant group scheme \(S_{N,\mathbb Z}\) is finite
etale even at these primes; it is the failure of linear
reductivity, not a failure of etaleness, that matters.
Computing a kernel dimension is an exact nonnegative linear-algebra
description.  Its ambient lattice can be very large, and this
description alone is not a positive combinatorial formula of the
kind required in the GCT flip.

\subsection*{Finding M-005: pure Artin realizations still retain a negative trace}
Let \(k=\mathbb F_q\), \(X=\operatorname{Spec}\mathbb F_{q^2}\), and
\(\ell\ne\operatorname{char}k\).
On the two geometric points of \(X\), geometric Frobenius is a
transposition.  The permutation representation
\(V=H^0(X_{\overline k},\mathbb Q_\ell)\) has eigenvalues \(1,-1\);
therefore \(\Lambda^2V\) is one-dimensional, pure of weight zero,
and has Frobenius trace \(-1\).

\begin{proof}
Frobenius exchanges the two basis vectors of \(V\).
In the basis consisting of their sum and difference, its matrix
is \(\operatorname{diag}(1,-1)\).
On \(\Lambda^2V\) it consequently acts by its determinant, \(-1\).
Every complex absolute value of this eigenvalue is one, which is
exactly purity of weight zero.  The assertions follow directly.
\end{proof}

In the category of rational Artin motives the same argument is the
decomposition \(h^0(X)=\mathbf1\oplus\varepsilon\), with
\(\Lambda^2h^0(X)=\varepsilon\); the rational idempotents are
\((1\pm \mathrm{Frob})/2\) in the finite quotient \(C_2\).
This is the motivic form of the corrected two-element Burnside
example, not a conjectural statement about general motives.
Purity controls absolute values and does not turn exterior powers
into effective finite sets or force nonnegative traces.  The RH
and Dirichlet-GRH input stipulated in the prompt supplies no
additional step in any proof in this paper.

\section{Dictionary into Grothendieck's frameworks}

\paragraph{Finite sets, group schemes and classifying topoi.}
For a finite group \(G\), SGA~3, Exp.~I, 4.1 (original p.~23) constructs its constant
group scheme \(G_{\mathbb Z}\).  Exp.~I, Proposition 4.7.2 (original pp.~34--35) identifies
representations of affine group schemes with comodules.
Thus a finite \(G\)-set \(E\) gives the locally free permutation
sheaf \(\mathcal O[E]\) on the classifying topos \(BG\).
Disjoint union becomes direct sum and Cartesian product becomes
tensor product.  Linearization is the resulting map into the
Grothendieck group of representations.  With short-exact-sequence
relations this is \(K_0(\operatorname{Rep}(G))\);
it must not be confused with the split group of indecomposable
representations in modular characteristic.

\paragraph{Exterior powers, symmetric powers and their signs.}
SGA~6, Exp.~V, Definition 2.1 (p.~314), Example 2.2(b), and Lemma 2.2.1 (pp.~315--316)
construct the pre-lambda structure using exterior powers and
their exact-sequence filtration.  Exp.~VI, Theorem 3.3
(printed pp.~392--393) proves the special lambda identities for
the \(K\)-group of locally free modules of bounded rank on a
commutatively ringed topos.  The exterior basis associated with
an \(r\)-element subset has the orientation sign of the induced
permutation of that subset.  Forgetting that sign gives the
incorrect subset formula audited in cote~133, scan~019.
The correct opposite-symmetric Burnside operations linearize to
these exterior powers, while remaining non-special upstairs.
This dictionary carries the proofs of A-001--A-002 and A-003 and the
signed formula M-001 to sheaves on \(BS_N\).

\paragraph{Schur functors and tensor powers.}
In characteristic zero, or after inverting \(N!\) for degree \(N\),
the rational symmetric-group projectors on a tensor power give
the usual Schur isotypic summands, compatibly with base change.
Over \(\mathbb Z\) one may instead use integral Schur modules and
filtrations; a rational direct-sum decomposition is not silently
asserted integrally.  Cote~133, scan~028 is the archive locator for
the tensor-power class and its monomial-symmetric expansion.
It concerns the induction product across the tower of symmetric
groups.  The Kronecker product is the tensor product at a fixed
degree, equivalently restriction along the diagonal
\(S_N\to S_N\times S_N\); cote~2, scan~017 records that
external-product/diagonal mechanism.

\paragraph{Gamma filtration.}
SGA~6, Exp.~V, 3.10 (p.~331) defines the gamma filtration and
3.10.1--3.10.2 give its generators.  Formula M-003 shows exactly
what happens to this filtration on \(K_0(\operatorname{Rep}_{\mathbb C}G)\)
after rationalization.  This calculation has no implication that
the integral filtration, or the gamma filtration of an arbitrary
scheme, vanishes.


\paragraph{Borger's lambda geometry.}
The representation ring, with its special lambda structure, is torsion-free
as an abelian group. Its Adams operations give commuting Frobenius lifts
on its spectrum: on characters, \(\psi^p(\chi)(g)=\chi(g^p)\).
This is the flat affine case of Borger, \emph{Lambda-rings and the field
with one element}, \S1.1 (the supplied source, lines 774--797).
The non-special Burnside structure audited in A-001--A-002 does not qualify for
that conclusion by the same route. Nor does the Frobenius-lift
interpretation restore effectivity: on the standard representation
\(V\) of \(S_3\), one has
\(\psi^2(V)=1+V-\operatorname{sgn}\), as computed in A-004.
Thus a scheme endomorphism arising from an Adams operation can have
negative coordinates in the irreducible-representation basis.
This is a concrete compatibility statement and a concrete limitation,
not a claim that all methods of lambda geometry fail.

\paragraph{Orbit-closure schemes and parameter spaces.}
The determinant and padded permanent are homogeneous polynomials
with integral coefficients.  Their rational orbit closures admit
integral models by contracting their defining ideals from
\(\mathbb Q\) to the ambient polynomial algebra over \(\mathbb Z\).
Those quotients are torsion-free as \(\mathbb Z\)-modules and hence
flat over \(\mathbb Z\); their generic fibers are the prescribed
closures.  It is not asserted that every special fiber is the
orbit closure of the reduced polynomial.
Projectivizing these homogeneous models places flat families with
fixed Hilbert polynomial in a Hilbert scheme.
The precise FGA source is Grothendieck, Bourbaki exposé~221,
Theorems 3.1 and 3.2 (printed pp.~259--260), for the Quot functor
and its fixed-Hilbert-polynomial projective pieces; the Hilbert
case is \(F=\mathcal O_X\), explained immediately after the
construction.  These representability results organize possible
families.  They do not prove the required comparison of the
determinant and permanent isotypic multiplicities.

\paragraph{Sheaves, weights and motives.}
The permutation sheaves and their orientation summands have
finite monodromy and admit Artin realizations.  Finding M-005
exhibits the sign already at this level.  For Grassmannian
realizations of the rectangular/two-row formula, Schubert cells
give even Tate cohomology and hard Lefschetz identifies the
partition-number difference with a primitive multiplicity.
The exact cohomological locators used here are SGA~4,
Exp.~XVIII, Theorem 3.2.5 (duality), SGA~5, Exp.~III, 4.7
(Lefschetz--Verdier trace formula), and Deligne,
\emph{La conjecture de Weil II}, 1.2.2 (purity) and
Theorem 4.1.1 (hard Lefschetz).
Their use does not identify an arbitrary Kronecker coefficient
with cohomology concentrated in one degree.
Only Artin motives and Tate realizations are used here. The
Grassmannian primitive quotient is a statement about cohomological
dimensions; no canonical motivic identification with a Kronecker
multiplicity space is asserted.
